% !TEX TS-program = xelatex
% !TEX encoding = UTF-8 Unicode
% !Mode:: "TeX:UTF-8"

\documentclass{resume}
\usepackage{zh_CN-Adobefonts_external} % Simplified Chinese Support using external fonts (./fonts/zh_CN-Adobe/)
% \usepackage{NotoSansSC_external}
% \usepackage{NotoSerifCJKsc_external}
% \usepackage{zh_CN-Adobefonts_internal} % Simplified Chinese Support using system fonts
\usepackage{linespacing_fix} % disable extra space before next section
\usepackage{cite}
\usepackage{multicol}

\begin{document}
\pagenumbering{gobble} % suppress displaying page number

\name{于海鑫}

\basicInfo{
  \email{name1e5s@bupt.edu.cn} \textperiodcentered\
  \phone{151-0113-5718} \textperiodcentered\
  \github[name1e5s]{https://github.com/name1e5s} \textperiodcentered\
  \homepage[blog.name1e5s.com]{http://blog.name1e5s.com/}}

\section{\faGraduationCap\  教育经历}
\datedsubsection{\textbf{北京邮电大学}, 北京}{2017 年 9 月 -- 2021 年 6 月}
\textit{本科生}\ 计算机学院，计算机科学与技术专业，GPA 3.50

\begin{onehalfspacing}
  主要课程成绩：
\end{onehalfspacing}
\begin{itemize}
  \item \textbf{大一} \quad 计算导论与程序设计（\textbf{91}）、
        计算导论与程序设计课程设计（\textbf{99}）、离散数学上（\textbf{95}）
  \item \textbf{大二} \quad 计算机系统基础（\textbf{92}）、计算机系统基础实践（\textbf{98}）、
        数据结构课程设计（\textbf{95}）、计算机组成原理课程设计（\textbf{97}）、
        计算机网络课程设计（\textbf{95}）、形式语言与自动机（\textbf{95}）、
        计算机组成原理（\textbf{90}）、计算机网络（\textbf{92}）
  \item \textbf{大三} \quad 数据库系统原理(\textbf{90})、算法设计与分析(\textbf{90})、
  Python 程序设计（\textbf{91}）、量子计算导论（\textbf{94}）、移动互联网技术及应用	（\textbf{91}）
\end{itemize}

\section{\faUsers\ 实习经历}
\datedsubsection{\textbf{清华大学信息技术研究院}}{2019 年 10 月 -- 2020 年 5 月}
\role{网络安全实验室}{实习生}
\begin{onehalfspacing}
  主要任务为探索利用 Intel SGX 技术对分布式图数据库 \emph{nebula} 做加固
\end{onehalfspacing}
\begin{itemize}
  \item 利用 AVX-2 指令集对 SGX 下的联邦学习进行加速
  \item 协助设计 SGX 场景下数据库的存储结构
\end{itemize}
\section{\faCode\ 个人项目}
\datedsubsection{\textbf{MIPS 微处理器}}{2019 年 5 月 -- 2019 年 8 月}
\role{SystemVerilog}{}
\begin{onehalfspacing}
  运行于 FPGA 上的顺序双发射 MIPS 微处理器, https://github.com/name1e5s/Sirius
  \begin{itemize}
    \item “龙芯杯”全国计算机系统能力培养大赛参赛作品，获\textbf{一等奖}（前三名）
    \item 从零开始自行设计的五级流水线双发射 MIPS32 微处理器，主频 100 MHz。
          同时基于该处理器搭建完整的 SoC，支持网口、串口等外设。可以在其上运行 uCore
          教学用操作系统，理论上能够运行完整的 Linux 操作系统
    \item 处理器性能较高，初赛性能分排名\textbf{第二}
  \end{itemize}
\end{onehalfspacing}

\datedsubsection{\textbf{DNS 中继服务器}}{2019 年 2 月 -- 2019 年 6 月}
\role{GoLang}{课程项目}
\begin{onehalfspacing}
  支持 DNS 转发及不良网址拦截的中继服务器, https://github.com/name1e5s/MuddyDNS
  \begin{itemize}
    \item “计算机网络课程设计” 课程项目
    \item 根据 RFC 1035 所述规范解析 DNS 包头，根据处理后的结果对 DNS 包进行操作
    \item 采用 CSP 并发模型，性能极高
  \end{itemize}
\end{onehalfspacing}

\datedsubsection{\textbf{车间调度模拟系统}}{2018 年 4 月 -- 2018 年 6 月}
\role{C, C++, Qt5}{课程项目}
\begin{onehalfspacing}
  支持求解 JobShop 问题的车间调度模拟系统, https://github.com/name1e5s/JobShop
  \begin{itemize}
    \item 使用 \textit{\href{http://www.jstor.org/stable/2632051}{The Shifting Bottleneck Procedure for Job Shop Scheduling}}
          所述的移动瓶颈法求解 JobShop 问题。
    \item 使用 Qt5 绘制简单界面以通过甘特图的形式展示调度结果
  \end{itemize}
\end{onehalfspacing}

\section{\faHeartO\ 获奖情况}
\datedline{“龙芯杯”第三届全国大学生计算机系统能力培养大赛\textbf{一等奖}（全国第二名）}{2019 年 8 月}
\datedline{校\textbf{二等奖学金}}{2019 年 9 月}
\datedline{校\textbf{三等奖学金}}{2018 年 9 月}
\section{\faCogs\ 个人技能}
% increase linespacing [parsep=0.5ex]
\begin{itemize}[parsep=0.5ex]
  \item 外国语言：CET-4\ \textbf{541} \quad CET-6\ \textbf{525}
  \item 编程语言:
        \begin{itemize}
          \item \textbf{熟悉}\quad C、Go、C++、Python、System Verilog
          \item \textbf{了解}\quad Swift、Chisel、Rust、Bash、TypeScript、Java
        \end{itemize}
  \item 平台/框架：Qt (C++)、Gin (Go)、Android App (Java)、iOS App (Swift)、FastAPI (Python)、Nest (TS)、FPGA (SV/Chisel)
  \item 其他：
        \begin{itemize}
          \item 熟悉 Linux 操作系统运行维护
          \item Linux 中国 (https://linux.cn) 译者之一，译作曾为该站 2018 年度阅读量\textbf{第一}
        \end{itemize}
\end{itemize}

%% Reference
%\newpage
%\bibliographystyle{IEEETran}
%\bibliography{mycite}
\end{document}
